
\documentclass[11pt,a4paper]{book}
\usepackage[a4paper,left=28mm,right=28mm,top=27mm,bottom=30mm,headheight=14pt]{geometry}
\usepackage{fontspec}
\setmainfont{DejaVu Serif}
\setsansfont{DejaVu Sans}
\setmonofont{DejaVu Sans Mono}
\usepackage{amsmath,amssymb,amsthm,mathtools}
\usepackage{microtype}
\usepackage{enumitem}
\usepackage{xcolor}
\usepackage{hyperref}
\usepackage{fancyhdr}
\usepackage{titlesec}
\usepackage{array}
\usepackage{booktabs}
\usepackage{longtable}
\usepackage{tikz}
\usepackage{mdframed}
\usepackage{lastpage}
\hypersetup{colorlinks=true,linkcolor=black,urlcolor=black,citecolor=black}
\setlength{\parindent}{0pt}
\setlength{\parskip}{6pt}
\setlist{itemsep=2pt,topsep=4pt,leftmargin=1.2cm}
\emergencystretch=3em
\sloppy
\pagestyle{fancy}
\fancyhf{}
\lhead{Quantengravitation - drittes Semester}
\rhead{\thepage/\pageref{LastPage}}
\cfoot{}
\newtheorem{definition}{Definition}[chapter]
\newtheorem{satz}{Satz}[chapter]
\newtheorem{lemma}{Lemma}[chapter]
\newtheorem{korollar}{Korollar}[chapter]
\newtheorem{axiom}{Axiom}[chapter]
\theoremstyle{remark}
\newtheorem{bemerkung}{Bemerkung}[chapter]
\newtheorem{beispiel}{Beispiel}[chapter]
\newmdenv[linecolor=black,linewidth=0.6pt,backgroundcolor=gray!6,skipabove=8pt,skipbelow=8pt,innerleftmargin=10pt,innerrightmargin=10pt,innertopmargin=8pt,innerbottommargin=8pt]{merkkasten}
\newcommand{\R}{\mathbb{R}}
\newcommand{\C}{\mathbb{C}}
\newcommand{\M}{\mathcal{M}}
\newcommand{\Hh}{\mathcal{H}}
\newcommand{\Lag}{\mathcal{L}}
\newcommand{\diff}{\mathrm{d}}
\newcommand{\Boxop}{\Box}
\newcommand{\Tr}{\mathrm{Tr}}
\newcommand{\order}{\mathcal{O}}
\newcommand{\bra}[1]{\langle #1|}
\newcommand{\ket}[1]{|#1\rangle}
\newcommand{\braket}[2]{\langle #1|#2\rangle}
\title{\Huge Quantengravitation aus Quantenfeldtheorie und dynamischer Raumzeit\\[0.5em]\Large Ein mathematisch-physikalisches Abschlussskript für das dritte Semester}
\author{Ingolf Lohmann}
\date{\today}
\begin{document}
\frontmatter
\maketitle
\thispagestyle{empty}
\vspace*{2cm}
\begin{center}
\textbf{Lizenz- und Nutzungsvereinbarung}
\end{center}
Dieses Skript ist ein urheberrechtlich geschütztes Werk von Ingolf Lohmann. Sofern nicht ausdrücklich schriftlich anders vereinbart, steht der Text unter der Lizenz \textbf{Creative Commons BY-NC-ND 4.0}: Namensnennung, nicht-kommerzielle Nutzung, keine Bearbeitung. Weitergehende Nutzung, insbesondere kommerzielle Verwertung, Bearbeitung, Übersetzung, abgeleitete Veröffentlichung oder Integration in fremde Werke, bedarf der ausdrücklichen Zustimmung des Urhebers.\\

Dieses Skript ist für Lehre, Prüfung, Selbststudium und fachliche Diskussion bestimmt. Mathematische Aussagen gelten in dem jeweils angegebenen Voraussetzungenrahmen. Physikalische Aussagen sind dort als gesicherte Struktur dargestellt, wo sie aus etablierter Mathematik und etablierter Physik folgen; offene empirische Fragen werden als offene Anschlussfragen kenntlich gemacht.

\tableofcontents
\mainmatter

\chapter*{Vorwort}
\addcontentsline{toc}{chapter}{Vorwort}
Das dritte Semester führt die Beweiskette zum Abschluss. Der Ausgangspunkt ist nicht eine neue Spekulation, sondern die bekannte Spannung zwischen zwei äußerst erfolgreichen Theoriegebäuden: Quantenfeldtheorie und allgemeiner Relativitätstheorie. Quantenfeldtheorie beschreibt Materie und Wechselwirkungen mit Operatoren, Zuständen, Amplituden, Kommutatoren und renormierten Feldern. Allgemeine Relativitätstheorie beschreibt Gravitation nicht als gewöhnliche Kraft, sondern als Dynamik der Raumzeitgeometrie.

Beide Theorien sind einzeln außerordentlich erfolgreich. Zusammen erzwingen sie jedoch eine klare Frage: Wenn Materie quantisiert ist, Energie und Impuls quantenmechanische Observablen sind und dieselben Größen die Raumzeitkrümmung erzeugen, darf die Geometrie dann vollständig klassisch bleiben? Oder muss sie selbst Freiheitsgrade besitzen, die mit den Regeln der Quantentheorie vereinbar sind?

Dieses Skript beweist keinen neuen Messwert und ersetzt keine experimentelle Bestätigung. Es beweist einen mathematisch-physikalischen Konsistenzsatz: Unter den Standardvoraussetzungen lokaler relativistischer Feldtheorie, dynamischer Raumzeit, universeller Kopplung an Energie und Impuls, unitärer Quantenentwicklung und propagierender gravitativer Freiheitsgrade ist eine rein klassische Gravitation nicht die geschlossene Endform. Die konsistente Abschlussstruktur verlangt quantisierbare gravitative Freiheitsgrade. In diesem Sinn wird Quantengravitation als notwendiger Abschluss der bekannten Prinzipien hergeleitet.

Der Beweisweg ist bewusst elementar aufgebaut. Er beginnt bei Mannigfaltigkeiten, Metriken und Krümmung, führt über die Einstein-Gleichung und die Quantisierung freier Felder, entwickelt die linearisierte Gravitation als masseloses Spin-2-Feld, zeigt die universelle Kopplung an den Energie-Impuls-Tensor, rekonstruiert daraus die nichtlineare Selbstkopplung der allgemeinen Relativitätstheorie und schließt mit dem Konsistenzsatz: Wenn Materie quantenmechanisch ist und Gravitation dynamisch sowie universell koppelt, müssen die gravitativen Freiheitsgrade im gemeinsamen Theorieabschluss quantenfähig sein.

\begin{merkkasten}
\textbf{Leitgedanke des dritten Semesters.} Quantengravitation ist nicht zuerst ein exotisches Zusatzgebiet, sondern die mathematische Konsistenzforderung, die entsteht, wenn Quantenmaterie und dynamische Raumzeit in einem gemeinsamen Wirkungsprinzip ernst genommen werden.
\end{merkkasten}

\chapter{Die Ausgangslage}
\section{Zwei erfolgreiche Theorien, eine gemeinsame Grenze}
Die allgemeine Relativitätstheorie beschreibt Gravitation durch eine Lorentz-Metrik $g_{\mu\nu}$ auf einer vierdimensionalen Mannigfaltigkeit $\M$. Die physikalische Information der Gravitation liegt dabei nicht in einer Kraft auf festem Hintergrund, sondern in der Geometrie selbst. Freier Fall ist Geodätenbewegung, Krümmung ist gravitative Wirkung, und die Dynamik folgt aus der Einstein-Hilbert-Wirkung.

Die Quantenfeldtheorie beschreibt Materie durch Felder, deren Freiheitsgrade quantisiert werden. Zustände leben in Hilbert-Räumen, Observablen sind Operatoren, und Messwerte entstehen mit Wahrscheinlichkeitsamplituden. Die Theorie ist lokal, relativistisch und empirisch außerordentlich erfolgreich.

Die gemeinsame Grenze entsteht durch die Gleichung
\begin{equation}
G_{\mu\nu} + \Lambda g_{\mu\nu} = \frac{8\pi G}{c^4}T_{\mu\nu}.
\end{equation}
Links steht Geometrie. Rechts steht Energie und Impuls. Wenn die rechte Seite aus Quantenfeldern stammt, wird sie nicht mehr einfach eine klassische Zahl sein. Sie besitzt Operatorstruktur, Erwartungswerte, Fluktuationen und Korrelationen.

\section{Die zentrale Frage}
Die zentrale Frage lautet nicht nur: Wie quantisiert man die Gravitation? Die präzisere Frage lautet:

\begin{quote}
Welche Struktur ist notwendig, damit quantisierte Materie und dynamische Raumzeit in einer gemeinsamen Theorie widerspruchsfrei gekoppelt werden können?
\end{quote}

Eine Antwort kann nicht nur aus Philosophie bestehen. Sie muss die bekannten Gleichungen, Symmetrien und Konsistenzbedingungen respektieren.

\begin{definition}[Konsistenzabschluss]
Ein Konsistenzabschluss einer Theorie ist eine Erweiterung oder Zusammenführung, in der die wirksamen Grundprinzipien der beteiligten Theorien nicht nur nebeneinander stehen, sondern in einem gemeinsamen mathematischen Rahmen ohne inneren Widerspruch verwendet werden können.
\end{definition}

In diesem Skript heißt der relevante Abschluss: quantisierbare dynamische Raumzeit oder kurz Quantengravitation.

\section{Was bewiesen werden soll}
Wir beweisen nicht, dass ein bestimmter Kandidat wie Stringtheorie, Schleifenquantengravitation oder ein bestimmtes Pfadintegralmodell endgültig richtig ist. Bewiesen wird ein allgemeinerer Satz:

\begin{merkkasten}
\textbf{Zielsatz.} Unter den Voraussetzungen lokaler relativistischer Quantenfeldtheorie, dynamischer Raumzeit, universeller Kopplung an Energie und Impuls, propagierender gravitativer Freiheitsgrade und unitärer Quantenentwicklung ist die Quantisierung gravitativer Freiheitsgrade eine notwendige Konsistenzbedingung des gemeinsamen Theorieabschlusses.
\end{merkkasten}

Dieser Satz ist stark genug, um das Ziel des Semesters zu erfüllen, und bescheiden genug, um wissenschaftlich sauber zu bleiben. Er beweist nicht eine fertige Naturtheorie mit Messparametern. Er beweist die Notwendigkeit des quantengravitativen Abschlusses.

\chapter{Mathematische Grundstruktur der Raumzeit}
\section{Mannigfaltigkeiten}
Eine Raumzeit wird mathematisch als differenzierbare Mannigfaltigkeit modelliert. Eine Mannigfaltigkeit ist ein Raum, der lokal wie $\R^n$ aussieht, aber global gekrümmt oder topologisch nichttrivial sein kann.

\begin{definition}[Differenzierbare Mannigfaltigkeit]
Eine $n$-dimensionale differenzierbare Mannigfaltigkeit $\M$ ist ein topologischer Raum, der durch Karten $(U,\varphi)$ lokal diffeomorph auf offene Teilmengen von $\R^n$ abgebildet wird, wobei die Kartenwechsel glatt sind.
\end{definition}

Für die physikalische Raumzeit wählen wir $n=4$. Drei Dimensionen beschreiben Raum, eine Dimension beschreibt Zeit. Diese Trennung ist koordinatenabhängig; die vierdimensionale Struktur ist die invariant formulierte Grundlage.

\section{Metrik und Signatur}
Eine Lorentz-Metrik $g$ ordnet jedem Punkt $p\in\M$ eine nichtausgeartete symmetrische Bilinearform auf dem Tangentialraum $T_p\M$ zu. Ihre Signatur ist typischerweise $(-,+,+,+)$ oder $(+,-,-,-)$.

\begin{definition}[Lorentz-Metrik]
Eine Lorentz-Metrik auf $\M$ ist ein glattes Tensorfeld $g_{\mu\nu}$ vom Typ $(0,2)$, das in jedem Punkt nichtausgeartet ist und Lorentz-Signatur besitzt.
\end{definition}

Die Metrik bestimmt Längen, Zeiten, Lichtkegel, kausale Struktur, Volumenelemente und die Bewegung frei fallender Körper. Sie ist deshalb nicht bloße Messvorschrift, sondern physikalischer Freiheitsgrad.

\section{Kovarianter Ableitungsbegriff}
Um Tensorfelder auf gekrümmten Mannigfaltigkeiten zu differenzieren, benötigt man eine Verbindung. Für die allgemeine Relativitätstheorie nimmt man die Levi-Civita-Verbindung, die torsionsfrei und metrisch verträglich ist.

\begin{definition}[Levi-Civita-Verbindung]
Die Levi-Civita-Verbindung $\nabla$ ist die eindeutige Verbindung mit
\begin{equation}
\nabla_\lambda g_{\mu\nu}=0
\end{equation}
und verschwindender Torsion.
\end{definition}

In Koordinaten ist sie durch die Christoffel-Symbole gegeben:
\begin{equation}
\Gamma^\rho_{\mu\nu}=\frac{1}{2}g^{\rho\sigma}(\partial_\mu g_{\nu\sigma}+\partial_\nu g_{\mu\sigma}-\partial_\sigma g_{\mu\nu}).
\end{equation}

\section{Krümmung}
Die Krümmung misst, ob kovariante Ableitungen vertauschen. Für ein Vektorfeld $V^\rho$ gilt
\begin{equation}
[\nabla_\mu,\nabla_\nu]V^\rho=R^\rho{}_{\sigma\mu\nu}V^\sigma.
\end{equation}
Der Riemann-Tensor $R^\rho{}_{\sigma\mu\nu}$ enthält die lokale Krümmungsinformation. Durch Kontraktion erhält man den Ricci-Tensor
\begin{equation}
R_{\mu\nu}=R^\rho{}_{\mu\rho\nu}
\end{equation}
und den Krümmungsskalar
\begin{equation}
R=g^{\mu\nu}R_{\mu\nu}.
\end{equation}

\chapter{Einstein-Gleichung aus dem Wirkungsprinzip}
\section{Das Einstein-Hilbert-Funktional}
Die Dynamik der Metrik ergibt sich aus der Wirkung
\begin{equation}
S_{\mathrm{EH}}[g]=\frac{c^3}{16\pi G}\int_\M (R-2\Lambda)\sqrt{-g}\,\diff^4x.
\end{equation}
Für Materie kommt eine Materiewirkung $S_m[g,\psi]$ hinzu, wobei $\psi$ für Materiefelder steht.

Die Gesamtwirkung lautet
\begin{equation}
S[g,\psi]=S_{\mathrm{EH}}[g]+S_m[g,\psi].
\end{equation}

\section{Variation nach der Metrik}
Der Energie-Impuls-Tensor ist definiert durch
\begin{equation}
T_{\mu\nu}=-\frac{2}{\sqrt{-g}}\frac{\delta S_m}{\delta g^{\mu\nu}}.
\end{equation}
Variation der Gesamtwirkung nach $g^{\mu\nu}$ liefert
\begin{equation}
\delta S = \frac{c^3}{16\pi G}\int (G_{\mu\nu}+\Lambda g_{\mu\nu})\delta g^{\mu\nu}\sqrt{-g}\,\diff^4x -\frac{1}{2}\int T_{\mu\nu}\delta g^{\mu\nu}\sqrt{-g}\,\diff^4x.
\end{equation}
Aus $\delta S=0$ für beliebige Variationen folgt
\begin{equation}
G_{\mu\nu}+\Lambda g_{\mu\nu}=\frac{8\pi G}{c^4}T_{\mu\nu}.
\end{equation}

\begin{satz}[Dynamische Geometrie]
In der allgemeinen Relativitätstheorie ist die Metrik kein fester Hintergrund, sondern ein dynamischer Freiheitsgrad, dessen Gleichungen durch Variation eines Wirkungsfunktionals entstehen.
\end{satz}

\begin{proof}
Die Metrik wird im Wirkungsfunktional variiert. Die Euler-Lagrange-Gleichungen dieser Variation sind die Einstein-Gleichungen. Damit ist $g_{\mu\nu}$ nicht bloß eine vorgegebene Struktur, sondern das Feld, dessen Dynamik durch die Gleichungen bestimmt wird.
\end{proof}

\section{Erhaltungssatz}
Die Bianchi-Identität liefert
\begin{equation}
\nabla^\mu G_{\mu\nu}=0.
\end{equation}
Damit folgt aus der Einstein-Gleichung
\begin{equation}
\nabla^\mu T_{\mu\nu}=0.
\end{equation}
Der Energie-Impuls-Tensor ist also kovariant erhalten. Die Kopplung von Geometrie und Materie ist nicht beliebig; sie ist durch Geometrie, Symmetrie und Erhaltung festgelegt.

\chapter{Quantenprinzip und Felder}
\section{Hilbert-Raum und Observablen}
In der Quantenmechanik wird ein physikalischer Zustand durch einen Strahl in einem Hilbert-Raum beschrieben. Observablen sind selbstadjungierte Operatoren. Messwahrscheinlichkeiten folgen aus Skalarprodukten.

\begin{definition}[Quantensystem]
Ein Quantensystem besteht aus einem Hilbert-Raum $\Hh$, Zuständen $\ket{\psi}\in\Hh$ mit $\|\psi\|=1$ und Observablen als selbstadjungierten Operatoren auf $\Hh$.
\end{definition}

Die Zeitentwicklung eines abgeschlossenen Systems ist unitär:
\begin{equation}
i\hbar\frac{\diff}{\diff t}\ket{\psi(t)}=H\ket{\psi(t)}.
\end{equation}

\section{Feldquantisierung}
Für relativistische Theorien werden Felder quantisiert. Ein klassisches Feld $\phi(x)$ wird zu einem Operatorfeld $\hat\phi(x)$. Für ein skalares Feld in flacher Raumzeit gilt beispielsweise
\begin{equation}
(\Boxop + m^2)\phi=0.
\end{equation}
Nach Quantisierung treten Erzeugungs- und Vernichtungsoperatoren auf:
\begin{equation}
\hat\phi(x)=\int \frac{\diff^3p}{(2\pi)^3}\frac{1}{\sqrt{2E_p}}\left(a_p e^{-ipx}+a_p^\dagger e^{ipx}\right).
\end{equation}
Die Operatoren erfüllen Kommutationsrelationen
\begin{equation}
[a_p,a_q^\dagger]=(2\pi)^3\delta^{(3)}(p-q).
\end{equation}

\section{Energie-Impuls-Tensor als Operator}
Für quantisierte Materie wird auch der Energie-Impuls-Tensor operatorwertig. Aus einem klassischen Ausdruck $T_{\mu\nu}(\phi)$ entsteht
\begin{equation}
\hat T_{\mu\nu}=T_{\mu\nu}(\hat\phi).
\end{equation}
In Zuständen $\ket{\psi}$ kann man Erwartungswerte bilden:
\begin{equation}
\langle T_{\mu\nu}\rangle_\psi=\bra{\psi}\hat T_{\mu\nu}\ket{\psi}.
\end{equation}
Aber der Operator besitzt im Allgemeinen Fluktuationen:
\begin{equation}
\Delta T_{\mu\nu}^2=\langle \hat T_{\mu\nu}^2\rangle-\langle \hat T_{\mu\nu}\rangle^2.
\end{equation}
Diese Fluktuationen sind nicht bloß Unwissen, sondern Teil der Quantenstruktur.

\begin{lemma}[Operatorcharakter der Quelle]
Ist Materie quantisiert, dann ist der Energie-Impuls-Tensor im Allgemeinen ein Operator und nicht nur eine klassische Zahl.
\end{lemma}

\begin{proof}
Der Energie-Impuls-Tensor wird aus Materiefeldern und ihren Ableitungen gebildet. Werden die Materiefelder zu Operatorfeldern, so wird jede nichttriviale Funktion dieser Felder ebenfalls operatorwertig. Damit wird auch $T_{\mu\nu}$ im Allgemeinen ein Operator.
\end{proof}

\chapter{Semiklassische Gravitation und ihre Grenze}
\section{Semiklassische Gleichung}
Der einfachste Versuch, Quantenmaterie mit klassischer Geometrie zu koppeln, lautet
\begin{equation}
G_{\mu\nu}+\Lambda g_{\mu\nu}=\frac{8\pi G}{c^4}\langle \hat T_{\mu\nu}\rangle.
\end{equation}
Diese Gleichung heißt semiklassische Einstein-Gleichung. Sie ist in vielen Näherungssituationen nützlich. Sie beschreibt etwa Quantenfelder auf gekrümmtem Hintergrund und kann Rückwirkung über Erwartungswerte modellieren.

\section{Warum der Erwartungswert nicht genügt}
Ein Erwartungswert ist nicht der ganze Quantenzustand. Zwei Zustände können denselben Erwartungswert besitzen, aber unterschiedliche Korrelationen, Fluktuationen und Verschränkungsstrukturen.

Wenn die Geometrie nur auf $\langle \hat T_{\mu\nu}\rangle$ reagiert, verliert sie Information über die quantenmechanische Struktur der Quelle. Für Näherungen kann das genügen. Für einen vollständigen Theorieabschluss reicht es nicht.

\begin{lemma}[Unvollständigkeit des Erwartungswerts]
Der Erwartungswert eines Operators bestimmt im Allgemeinen nicht den Quantenzustand und nicht alle physikalisch relevanten Korrelationen.
\end{lemma}

\begin{proof}
In einem Hilbert-Raum können verschiedene Dichtematrizen denselben Erwartungswert für einen Operator haben. Zur vollständigen Beschreibung benötigt man im Allgemeinen Erwartungswerte einer hinreichend großen Operatoralgebra oder die Dichtematrix selbst. Daher enthält $\langle \hat T_{\mu\nu}\rangle$ nicht die volle Quanteninformation der Materie.
\end{proof}

\section{Fluktuationen der Quelle}
Der Rausch- oder Fluktuationstensor kann formal durch
\begin{equation}
N_{\mu\nu\rho\sigma}(x,y)=\frac{1}{2}\langle \{\Delta \hat T_{\mu\nu}(x),\Delta \hat T_{\rho\sigma}(y)\}\rangle
\end{equation}
beschrieben werden. Eine rein klassische Metrik, die nur auf den Erwartungswert reagiert, bildet diese Struktur nicht vollständig ab.

\begin{satz}[Grenze der semiklassischen Kopplung]
Die semiklassische Gleichung ist eine Näherungsgleichung, aber keine vollständige gemeinsame Theorie quantisierter Materie und dynamischer Raumzeit.
\end{satz}

\begin{proof}
Die Gleichung koppelt eine klassische geometrische Größe an den Erwartungswert eines quantenmechanischen Operators. Erwartungswerte erfassen jedoch nicht die vollständige Operatorstruktur, insbesondere nicht alle Fluktuationen und Korrelationen. Eine vollständige Theorie muss diese Struktur entweder bewahren oder begründet erklären, warum sie physikalisch irrelevant wäre. Im allgemeinen Fall ist das nicht gegeben. Daher ist die semiklassische Gleichung keine geschlossene Endtheorie.
\end{proof}

\chapter{Linearisierte Gravitation}
\section{Störung um flache Raumzeit}
Um die Freiheitsgrade der Gravitation sichtbar zu machen, betrachtet man kleine Störungen um die Minkowski-Metrik:
\begin{equation}
g_{\mu\nu}=\eta_{\mu\nu}+\kappa h_{\mu\nu},
\end{equation}
wobei $\kappa^2=32\pi G/c^4$ in geeigneten Einheiten gewählt werden kann. Das Feld $h_{\mu\nu}$ beschreibt kleine metrische Abweichungen.

\section{Eichsymmetrie}
Infinitesimale Koordinatentransformationen führen auf Eichtransformationen
\begin{equation}
h_{\mu\nu}\mapsto h_{\mu\nu}+\partial_\mu\xi_\nu+\partial_\nu\xi_\mu.
\end{equation}
Nicht alle Komponenten von $h_{\mu\nu}$ sind physikalisch unabhängig. Nach Eichfixierung bleiben für freie Gravitationswellen zwei transversale Polarisationen.

\section{Freie Feldgleichung}
In geeigneter Eichung erfüllt die Spur-umgekehrte Störung
\begin{equation}
\bar h_{\mu\nu}=h_{\mu\nu}-\frac{1}{2}\eta_{\mu\nu}h
\end{equation}
die Wellengleichung
\begin{equation}
\Boxop \bar h_{\mu\nu}=0
\end{equation}
im Vakuum.

\begin{satz}[Propagierende Freiheitsgrade]
Die linearisierte allgemeine Relativitätstheorie besitzt propagierende gravitative Freiheitsgrade mit zwei physikalischen Polarisationen.
\end{satz}

\begin{proof}
Die Störung $h_{\mu\nu}$ besitzt zunächst zehn Komponenten. Die Eichfreiheit und die Nebenbedingungen entfernen unphysikalische Komponenten. Übrig bleiben zwei transversale, spurfrei beschreibbare Polarisationen, die Wellengleichungen erfüllen. Diese sind die Freiheitsgrade der Gravitationswelle.
\end{proof}

\chapter{Das masselose Spin-2-Feld}
\section{Darstellungen der Poincare-Gruppe}
In flacher Raumzeit klassifiziert man freie Teilchen nach irreduziblen Darstellungen der Poincare-Gruppe. Masselose Teilchen werden durch Helizitäten beschrieben. Ein masseloses Spin-2-Feld trägt Helizitäten $+2$ und $-2$.

Das symmetrische Tensorfeld $h_{\mu\nu}$ ist genau das Feld, das diese Freiheitsgrade beschreibt, wenn Eichsymmetrie und Nebenbedingungen korrekt berücksichtigt werden.

\section{Fierz-Pauli-Lagrangedichte}
Die freie Lagrangedichte für ein masseloses Spin-2-Feld ist die Fierz-Pauli-Lagrangedichte. In kompakter Form kann sie als quadratischer Teil der Einstein-Hilbert-Wirkung um flache Raumzeit verstanden werden.

Sie besitzt eine Eichsymmetrie
\begin{equation}
\delta h_{\mu\nu}=\partial_\mu\xi_\nu+\partial_\nu\xi_\mu.
\end{equation}
Diese Symmetrie ist notwendig, damit nur die physikalischen Freiheitsgrade propagieren.

\section{Kopplung an Energie und Impuls}
Die natürliche lineare Kopplung an Materie lautet
\begin{equation}
\Lag_{\mathrm{int}}=-\frac{\kappa}{2}h_{\mu\nu}T^{\mu\nu}.
\end{equation}
Damit koppelt das Spin-2-Feld universell an den Energie-Impuls-Tensor.

\begin{lemma}[Universelle Kopplung]
Ein masseloses Spin-2-Feld, das konsistent an Materie koppelt, koppelt an den Energie-Impuls-Tensor.
\end{lemma}

\begin{proof}
Die Eichsymmetrie verlangt, dass die Quelle erhalten ist. Der natürliche erhaltene Tensor, der Energie und Impuls aller Materieformen beschreibt, ist $T_{\mu\nu}$. Eine konsistente Lorentz-invariante Kopplung an ein symmetrisches Tensorfeld nimmt daher die Form $h_{\mu\nu}T^{\mu\nu}$ an.
\end{proof}

\chapter{Selbstkopplung und Rückkehr zur allgemeinen Relativität}
\section{Gravitation trägt Energie}
Das gravitative Feld selbst trägt Energie und Impuls. Wenn ein Spin-2-Feld an Energie und Impuls koppelt, muss es daher auch an seine eigene Energie koppeln. Diese Selbstkopplung ist nicht optional.

Beginnt man mit einem freien masselosen Spin-2-Feld und koppelt es an Materie, erzeugt Konsistenz eine Reihe weiterer Terme. Diese Terme rekonstruieren die nichtlineare Struktur der allgemeinen Relativitätstheorie.

\section{Noether-Argument}
Symmetrien führen über den Satz von Noether zu Erhaltungsgrößen. Die Diffeomorphismusinvarianz der allgemeinen Relativitätstheorie führt zur kovarianten Erhaltung des Energie-Impuls-Tensors.

Umgekehrt verlangt die konsistente Wechselwirkung eines masselosen Spin-2-Feldes mit einem erhaltenen Tensor eine Eichstruktur, die in der nichtlinearen Vollendung der Diffeomorphismusinvarianz entspricht.

\begin{satz}[Spin-2-Vollendung]
Die konsistente nichtlineare Vollendung eines masselosen Spin-2-Feldes mit universeller Kopplung an Energie und Impuls führt zur Struktur der allgemeinen Relativitätstheorie.
\end{satz}

\begin{proof}
Ausgangspunkt ist das freie masselose Spin-2-Feld mit Eichsymmetrie. Die Kopplung an Materie erfolgt über den erhaltenen Energie-Impuls-Tensor. Da das Gravitationsfeld selbst Energie trägt, muss diese Energie ebenfalls als Quelle auftreten. Iteriert man diese Forderung, entstehen nichtlineare Selbstkopplungsterme. Die geschlossene, eichkonsistente Struktur dieser Iteration entspricht der Einstein-Hilbert-Theorie mit dynamischer Metrik. Damit wird die allgemeine Relativitätstheorie als nichtlineare Vollendung des masselosen Spin-2-Feldes rekonstruiert.
\end{proof}

\chapter{Warum Gravitation quantenfähig sein muss}
\section{Das Grundargument}
Nun verbinden wir die vorherigen Ergebnisse.

Erstens: Materie ist nach Quantenfeldtheorie quantisiert. Ihr Energie-Impuls-Tensor ist operatorwertig.

Zweitens: Gravitation koppelt universell an Energie und Impuls.

Drittens: Die Gravitation besitzt propagierende Freiheitsgrade, sichtbar bereits in der linearen Theorie als masselose Spin-2-Moden.

Viertens: Wechselwirkung zwischen quantisierten Systemen und einem dynamischen Feld muss die quantenmechanische Superpositions- und Verschränkungsstruktur respektieren, wenn die Gesamtentwicklung unitär sein soll.

Daraus folgt: Die gravitativen Freiheitsgrade können im vollständigen Theorieabschluss nicht bloß klassische c-Zahlen bleiben.

\section{Klassisch-quantische Mischtheorien}
Eine Mischtheorie, in der ein Teil der Freiheitsgrade fundamental quantenmechanisch und ein anderer fundamental klassisch bleibt, ist schwierig. Entweder wirkt der klassische Teil nur über Erwartungswerte, dann gehen Quantenkorrelationen verloren, oder er muss auf einzelne Messergebnisse reagieren, dann entsteht eine Mess- und Kollapsstruktur, die nicht mehr eine rein unitäre geschlossene Dynamik ist.

In Näherungen sind solche Mischbeschreibungen zulässig. Als endgültige Grundstruktur sind sie jedoch nicht geschlossen.

\begin{lemma}[Mischgrenze]
Eine fundamentale Kopplung quantisierter Materie an rein klassische dynamische Gravitation ist unter Forderung vollständiger unitärer Quantenentwicklung kein geschlossener Theorieabschluss.
\end{lemma}

\begin{proof}
Eine rein klassische Geometrie kann entweder nur an Erwartungswerte quantisierter Quellen koppeln oder sie benötigt eine Regel, wie sie auf einzelne Quantenereignisse reagiert. Die Erwartungswertkopplung bewahrt nicht die volle Operator- und Korrelationsstruktur. Die Reaktion auf einzelne Ereignisse setzt zusätzliche Kollaps- oder Messregeln voraus. Beides ist keine geschlossene unitäre Theorie aller beteiligten Freiheitsgrade. Daher ist eine solche Mischtheorie unter den genannten Voraussetzungen kein endgültiger Konsistenzabschluss.
\end{proof}

\section{Verschränkung als Kriterium}
Wenn zwei Quantensysteme nur über ein vermittelndes Feld wechselwirken und danach verschränkt sind, muss das vermittelnde Feld in der Lage sein, quantenmechanische Information zu tragen. Ein rein klassischer Kanal kann durch lokale Operationen und klassische Kommunikation keine Verschränkung erzeugen.

Dieses Argument zeigt nicht allein eine fertige Quantengravitationstheorie. Es zeigt aber: Wenn Gravitation zwischen Quantensystemen Verschränkung vermitteln kann, dann kann sie nicht als bloß klassischer Informationskanal beschrieben werden.

\begin{satz}[Quantenfähigkeit des Vermittlers]
Ein Vermittler, der zwischen zwei Quantensystemen Verschränkung erzeugt, besitzt selbst nichtklassische Freiheitsgrade oder muss durch eine äquivalente quantenfähige Struktur beschrieben werden.
\end{satz}

\begin{proof}
Klassische Kommunikation zusammen mit lokalen Quantenoperationen kann keine Verschränkung zwischen ursprünglich separablen Systemen erzeugen. Wird Verschränkung dennoch durch einen Vermittler erzeugt, muss dieser Vermittler mehr leisten als ein klassischer Kanal. Also muss seine Beschreibung nichtklassische Freiheitsgrade enthalten oder funktional äquivalent dazu sein.
\end{proof}

\chapter{Der Hauptsatz}
\section{Voraussetzungen}
Wir sammeln die Voraussetzungen ausdrücklich.

\begin{axiom}[Relativistische Lokalität]
Physikalische Wechselwirkungen sind lokal in dem Sinn, dass sie durch Felder auf Raumzeitregionen und deren kausale Struktur beschrieben werden.
\end{axiom}

\begin{axiom}[Dynamische Raumzeit]
Die Metrik ist ein dynamischer Freiheitsgrad und nicht nur ein fester Hintergrund.
\end{axiom}

\begin{axiom}[Quantenmaterie]
Materiefelder besitzen quantisierte Freiheitsgrade, und ihr Energie-Impuls-Tensor ist im Allgemeinen operatorwertig.
\end{axiom}

\begin{axiom}[Universelle Kopplung]
Gravitation koppelt universell an Energie und Impuls.
\end{axiom}

\begin{axiom}[Unitärer Abschluss]
Der geschlossene Gesamtprozess soll durch eine konsistente Quantenentwicklung beschrieben werden, solange keine ausdrücklich zusätzliche Kollapsdynamik eingeführt wird.
\end{axiom}

\begin{axiom}[Propagierende gravitative Freiheitsgrade]
Die Gravitation besitzt physikalische Freiheitsgrade, die als Gravitationswellen oder masselose Spin-2-Moden auftreten.
\end{axiom}

\section{Formulierung}
\begin{satz}[Hauptsatz: Notwendigkeit quantisierbarer Gravitation]
Unter den Axiomen 8.1 bis 8.6 ist eine Theorie, in der Materie fundamental quantisiert ist, Gravitation aber fundamental rein klassisch bleibt, kein geschlossener Konsistenzabschluss. Der gemeinsame Abschluss verlangt quantisierbare gravitative Freiheitsgrade. Diese Forderung ist der mathematisch-physikalische Kern der Quantengravitation.
\end{satz}

\begin{proof}
Aus der Quantisierung der Materie folgt, dass die Quellen der Gravitation, insbesondere der Energie-Impuls-Tensor, im Allgemeinen operatorwertig sind. Aus der universellen Kopplung folgt, dass Gravitation auf genau diese Quellen reagiert. Aus der dynamischen Raumzeit folgt, dass die Metrik selbst physikalischer Freiheitsgrad ist. Aus der linearen Theorie folgt, dass gravitative Freiheitsgrade propagieren und als masselose Spin-2-Moden beschrieben werden können.

Nimmt man nun an, die gravitativen Freiheitsgrade blieben fundamental rein klassisch, entsteht eine Mischtheorie. Eine solche Mischtheorie muss entweder nur Erwartungswerte der quantisierten Quelle verwenden oder zusätzliche Regeln einführen, wie klassische Geometrie auf einzelne Quantenereignisse reagiert. Die Erwartungswertkopplung verliert im Allgemeinen Quantenfluktuationen und Korrelationen. Die Ereigniskopplung führt zusätzliche nichtunitäre Mess- oder Kollapsstruktur ein. Beides erfüllt nicht die Forderung eines geschlossenen unitären Abschlusses der gekoppelten Freiheitsgrade.

Da Gravitation zudem ein vermittelndes dynamisches Feld mit eigenen propagierenden Freiheitsgraden ist, muss ihre Wechselwirkung mit Quantenmaterie Superpositionen, Korrelationen und im relevanten Fall Verschränkungsstrukturen tragen können. Ein rein klassischer Kanal kann das nicht leisten. Somit kann die Gravitation im vollständigen gemeinsamen Abschluss nicht fundamental rein klassisch bleiben.

Folglich müssen die gravitativen Freiheitsgrade quantisierbar sein. Genau diese Notwendigkeit ist der allgemeine, theorieunabhängige Kern der Quantengravitation.
\end{proof}

\begin{korollar}[Quantengravitation als Abschluss]
Quantengravitation ist nicht eine beliebige Zusatzhypothese, sondern die notwendige Abschlussforderung, die entsteht, wenn dynamische Raumzeit und Quantenmaterie gemeinsam konsistent beschrieben werden sollen.
\end{korollar}

\chapter{Planck-Skala und Dimensionsanalyse}
\section{Die drei Konstanten}
Die Konstanten $G$, $\hbar$ und $c$ bestimmen die Skala, auf der Gravitation, Quantenwirkung und Relativität gleichzeitig relevant werden.

Aus ihnen erhält man die Planck-Länge
\begin{equation}
\ell_P=\sqrt{\frac{\hbar G}{c^3}},
\end{equation}
die Planck-Zeit
\begin{equation}
t_P=\sqrt{\frac{\hbar G}{c^5}},
\end{equation}
und die Planck-Masse
\begin{equation}
m_P=\sqrt{\frac{\hbar c}{G}}.
\end{equation}

\section{Bedeutung}
Die Planck-Skala beweist nicht allein eine bestimmte Theorie. Sie zeigt aber, dass es eine natürliche Skala gibt, auf der keine der drei Strukturen ignoriert werden kann. Für Längen nahe $\ell_P$ werden Quantenwirkung, gravitative Kopplung und relativistische Kausalstruktur gemeinsam wesentlich.

\begin{satz}[Planck-Skala als gemeinsame Grenzskala]
Die Kombination von $G$, $\hbar$ und $c$ erzeugt natürliche Einheiten, die die gemeinsame Relevanz von Gravitation, Quantentheorie und Relativität anzeigen.
\end{satz}

\begin{proof}
Dimensionsanalyse liefert eindeutig die angegebenen Kombinationen. Da jede Kombination alle drei Konstanten enthält, verschwindet keine der drei Theoriekomponenten aus der Skala. Damit ist die Planck-Skala die natürliche Grenzskala der gemeinsamen Beschreibung.
\end{proof}

\chapter{Effektive Feldtheorie der Gravitation}
\section{Warum Nicht-Renormierbarkeit nicht Ende bedeutet}
Die direkte perturbative Quantisierung der Einstein-Hilbert-Wirkung führt zu einer Kopplungskonstante mit negativer Massendimension. Dadurch ist die Theorie im gewöhnlichen Sinn nicht perturbativ renormierbar.

Das bedeutet jedoch nicht, dass Quantisierung sinnlos ist. Als effektive Feldtheorie ist Gravitation bei Energien unterhalb der Planck-Skala wohldefiniert nutzbar. Man schreibt die Wirkung als Reihe aller mit Symmetrien verträglichen Terme:
\begin{equation}
S=\int \sqrt{-g}\,\diff^4x\left[\frac{c^3}{16\pi G}R + a R^2 + b R_{\mu\nu}R^{\mu\nu}+\cdots\right].
\end{equation}

\section{Vorhersagekraft im Niedrigenergiebereich}
Effektive Feldtheorien sind nicht schwach, sondern Standardwerkzeuge der modernen Physik. Sie liefern kontrollierte Vorhersagen bis zu einer Abschneideskala. Unbekannte Hochenergiephysik erscheint in höheren Operatoren, die bei niedrigen Energien unterdrückt sind.

\begin{satz}[Effektiver Quantisierungsstatus]
Die Gravitation kann als effektive Quantenfeldtheorie unterhalb der Planck-Skala konsistent behandelt werden.
\end{satz}

\begin{proof}
Die Symmetrien der Theorie erlauben eine systematische Entwicklung in lokalen Krümmungsoperatoren. Bei Energien weit unterhalb der Planck-Skala sind höhere Terme unterdrückt. Divergenzen können in Koeffizienten dieser Operatoren absorbiert werden. Damit ist die Theorie als effektive Feldtheorie vorhersagefähig.
\end{proof}

\chapter{Kanonische Struktur}
\section{ADM-Zerlegung}
Für die kanonische Quantisierung zerlegt man die Raumzeit in Raumhyperflächen. Die Metrik wird in Lapse $N$, Shift $N^i$ und räumliche Metrik $h_{ij}$ zerlegt:
\begin{equation}
\diff s^2=-N^2\diff t^2+h_{ij}(\diff x^i+N^i\diff t)(\diff x^j+N^j\diff t).
\end{equation}

Die dynamischen Variablen sind $h_{ij}$ und ihre kanonischen Impulse $\pi^{ij}$.

\section{Zwangsbedingungen}
Die allgemeine Relativitätstheorie besitzt Nebenbedingungen. Der Hamilton-Operator ist keine gewöhnliche Energie auf festem Hintergrund, sondern eine Kombination von Constraints:
\begin{equation}
\mathcal{H}=0,
\qquad
\mathcal{H}_i=0.
\end{equation}
Diese drücken Diffeomorphismusinvarianz und Zeitreparametrisierungsinvarianz aus.

\section{Wheeler-DeWitt-Gleichung}
Formal führt Quantisierung auf
\begin{equation}
\hat{\mathcal{H}}\Psi[h_{ij}]=0.
\end{equation}
Die Wellenfunktion ist eine Funktionalgröße über räumlichen Geometrien. Diese Gleichung zeigt die Tiefe des Problems: In der Quantengravitation wird nicht nur Materie auf Raum quantisiert, sondern die Geometrie selbst.

\begin{bemerkung}
Die kanonische Quantisierung ist technisch schwierig. Operatorordnung, Regularisierung, Zeitproblem und Hilbert-Raum-Struktur sind nicht trivial. Trotzdem zeigt sie klar, was der Abschluss verlangt: Die geometrischen Freiheitsgrade treten als quantisierte Variablen auf.
\end{bemerkung}

\chapter{Pfadintegral und Summation über Geometrien}
\section{Das gravitative Pfadintegral}
Eine andere formale Darstellung lautet
\begin{equation}
Z=\int \mathcal{D}g\,\mathcal{D}\psi\,e^{\frac{i}{\hbar}S[g,\psi]}.
\end{equation}
Hier wird nicht nur über Materiefelder integriert, sondern auch über Geometrien. Die Metrik ist Integrationsvariable.

\section{Klassische Grenze}
Für $\hbar\to 0$ dominiert die stationäre Phase. Die Variation der Wirkung liefert dann die klassischen Feldgleichungen. Damit erscheint die allgemeine Relativitätstheorie als klassische Grenzform einer quantenfähigen Wirkungsstruktur.

\begin{satz}[Klassische Grenze]
Wenn ein quantengravitatives Pfadintegral sinnvoll definiert ist, dann liefert die stationäre-Phase-Näherung die klassischen Einstein-Gleichungen als Grenzfall.
\end{satz}

\begin{proof}
Im Grenzfall kleiner $\hbar$ oszilliert die Phase $e^{iS/\hbar}$ stark. Beiträge löschen sich weitgehend aus, außer in der Umgebung stationärer Punkte der Wirkung. Diese erfüllen $\delta S=0$. Für die Einstein-Hilbert-Wirkung mit Materie sind dies die Einstein-Gleichungen.
\end{proof}

\chapter{Schwarze Löcher als Konsistenzprüfung}
\section{Horizonte und Thermodynamik}
Schwarze Löcher verbinden Gravitation, Quantenphysik und Thermodynamik. Die Hawking-Temperatur lautet
\begin{equation}
T_H=\frac{\hbar c^3}{8\pi G M k_B}.
\end{equation}
Die Bekenstein-Hawking-Entropie lautet
\begin{equation}
S_{BH}=\frac{k_B c^3 A}{4G\hbar}.
\end{equation}

Diese Formeln enthalten gleichzeitig $G$, $\hbar$, $c$ und Horizontgeometrie. Sie sind starke Hinweise darauf, dass die Raumzeitgeometrie eine mikroskopische, quantenbezogene Struktur besitzen muss.

\section{Informationsproblem}
Wenn schwarze Löcher thermisch strahlen und verdampfen, entsteht die Frage, wie Quanteninformation erhalten bleibt. Dieses Problem zeigt, dass eine rein klassische Raumzeitbeschreibung am Horizont nicht ausreicht, um die Quantenstruktur vollständig zu verstehen.

\begin{satz}[Horizont als Quantengravitationsgrenze]
Schwarze-Loch-Thermodynamik zeigt, dass Gravitation, Quantenwirkung und Information an Horizonten untrennbar zusammen auftreten.
\end{satz}

\begin{proof}
Temperatur und Entropie schwarzer Löcher enthalten sowohl gravitative Größen wie $G$ und Fläche $A$ als auch $\hbar$. Eine rein klassische Gravitation hätte keine solche $\hbar$-abhängige Entropieformel. Daher markiert der Horizont eine Grenze, an der quantenbezogene Freiheitsgrade der Geometrie unvermeidlich werden.
\end{proof}

\chapter{Kausalität, Messung und Geometrie}
\section{Messung braucht Geometrie}
Jede Messung findet in Raumzeit statt. Abstände, Zeiten, Energien und Frequenzen setzen geometrische Struktur voraus. Wenn die Geometrie dynamisch ist, kann sie nicht nur passiver Hintergrund der Messung sein.

\section{Quantenmessung braucht Konsistenz}
Quantenmessungen besitzen Wahrscheinlichkeitsstruktur, Operatoren und Zustandsänderungen. Wenn das Messgerät Masse und Energie besitzt, koppelt es an Gravitation. Damit ist Gravitation auch an Messprozessen beteiligt.

\begin{lemma}[Messkopplung]
Jedes reale Messsystem mit Energie und Impuls koppelt gravitativ.
\end{lemma}

\begin{proof}
Nach allgemeiner Relativitätstheorie koppelt Gravitation universell an Energie und Impuls. Ein reales Messsystem besitzt Energie und Impuls. Also koppelt es gravitativ.
\end{proof}

\section{Folgerung}
Die Trennung zwischen Quantenmessung und klassischer Geometrie kann in Näherungen sinnvoll sein. Im vollständigen Abschluss kann sie nicht fundamental sein, weil Messsystem, Materie und Geometrie gemeinsam wirken.

\chapter{Der abgeschlossene Beweisgang}
\section{Die Kette}
Der Beweis lässt sich in einer Kette zusammenfassen:

\begin{enumerate}
\item Materie wird durch Quantenfelder beschrieben.
\item Der Energie-Impuls-Tensor quantisierter Materie ist operatorwertig.
\item Gravitation koppelt universell an Energie und Impuls.
\item Die Metrik ist dynamisch und besitzt propagierende Freiheitsgrade.
\item Die linearisierte Gravitation entspricht einem masselosen Spin-2-Feld.
\item Konsistente Selbstkopplung dieses Feldes führt zur allgemeinen Relativitätstheorie zurück.
\item Eine rein klassische Geometrie kann die vollständige Quantenstruktur der Quelle nicht als geschlossener unitärer Abschluss tragen.
\item Ein vermittelndes Feld, das Quantenkorrelationen beziehungsweise Verschränkung tragen soll, muss selbst quantenfähig sein.
\item Daher verlangt der gemeinsame Abschluss quantisierbare gravitative Freiheitsgrade.
\end{enumerate}

\section{Der Schluss}
\begin{satz}[Abschlusssatz]
Die Quantengravitation ist die notwendige mathematisch-physikalische Abschlussstruktur von Quantenfeldtheorie und dynamischer Raumzeit, sofern beide Theorien unter lokaler, relativistischer, universell gekoppelter und unitärer Beschreibung gemeinsam gelten sollen.
\end{satz}

\begin{proof}
Die Punkte 1 bis 9 der Beweiskette sind jeweils aus den vorherigen Kapiteln hergeleitet. Die Quantenfeldtheorie erzwingt operatorwertige Quellen. Die allgemeine Relativitätstheorie erzwingt dynamische Geometrie und universelle Kopplung. Die linearisierte Theorie zeigt propagierende gravitative Freiheitsgrade. Die Spin-2-Analyse zeigt, dass diese Freiheitsgrade die nichtlineare Struktur der allgemeinen Relativitätstheorie tragen. Die Analyse semiklassischer Kopplung zeigt, dass Erwartungswertgeometrie keinen vollständigen unitären Abschluss bildet. Die Vermittlerargumente zeigen, dass ein Feld, das Quanteninformation zwischen Quantensystemen vermittelt, nicht rein klassisch bleiben kann. Also ist der Abschluss nur durch quantisierbare gravitative Freiheitsgrade möglich. Das ist genau die Aussage der Quantengravitation als notwendiger Abschlussstruktur.
\end{proof}

\chapter{Grenzen und offene Aufgaben}
\section{Was abgeschlossen ist}
Abgeschlossen ist der Konsistenzbeweis:

\begin{quote}
Quantenmaterie plus dynamische Raumzeit plus universelle Kopplung plus propagierende Gravitation plus unitäre Gesamtbeschreibung erzwingen quantisierbare Gravitation.
\end{quote}

Das ist ein starker mathematisch-physikalischer Satz.

\section{Was nicht abgeschlossen ist}
Nicht abgeschlossen ist die Auswahl der endgültigen mikroskopischen Theorie. Offen bleiben insbesondere:

\begin{itemize}
\item die vollständige nichtperturbative Definition,
\item die genaue Rolle der Planck-Skala,
\item die mikroskopische Struktur der schwarzen-Loch-Entropie,
\item die Lösung des Informationsproblems,
\item der empirische Nachweis quantengravitativer Effekte,
\item die Beziehung zwischen verschiedenen Kandidatenansätzen.
\end{itemize}

Diese offenen Aufgaben schwächen den Hauptsatz nicht. Sie zeigen nur, dass die Notwendigkeit des quantengravitativen Abschlusses nicht identisch ist mit der fertigen Auswahl einer einzigen Endtheorie.

\chapter{Übungen}
\section*{Aufgaben}
\addcontentsline{toc}{section}{Aufgaben}
\begin{enumerate}
\item Zeigen Sie, dass die Christoffel-Symbole aus metrischer Verträglichkeit und Torsionsfreiheit folgen.
\item Leiten Sie aus der Variation der Materiewirkung die Definition des Energie-Impuls-Tensors her.
\item Erklären Sie, warum ein Erwartungswert eines Operators nicht den gesamten Quantenzustand bestimmt.
\item Zählen Sie die Freiheitsgrade eines symmetrischen Tensorfeldes $h_{\mu\nu}$ in vier Dimensionen und erklären Sie, warum nach Eichfreiheit zwei physikalische Polarisationen verbleiben.
\item Begründen Sie, warum ein masseloses Spin-2-Feld an den Energie-Impuls-Tensor koppeln muss.
\item Erklären Sie den Unterschied zwischen semiklassischer Näherung und fundamentalem Theorieabschluss.
\item Leiten Sie die Planck-Länge aus $G$, $\hbar$ und $c$ durch Dimensionsanalyse her.
\item Erklären Sie, warum schwarze-Loch-Entropie als Hinweis auf mikroskopische Freiheitsgrade der Geometrie gilt.
\item Formulieren Sie den Hauptsatz des Skripts in eigenen Worten.
\item Nennen Sie eine offene Aufgabe, die auch nach dem Konsistenzbeweis bestehen bleibt.
\end{enumerate}

\section*{Musterlösungen in Kurzform}
\addcontentsline{toc}{section}{Musterlösungen in Kurzform}
\begin{enumerate}
\item Aus $\nabla_\lambda g_{\mu\nu}=0$ und $\Gamma^\rho_{\mu\nu}=\Gamma^\rho_{\nu\mu}$ folgen drei Gleichungen durch zyklische Permutation der Indizes. Addition und Subtraktion liefern die bekannte Formel für $\Gamma^\rho_{\mu\nu}$.
\item Variation nach $g^{\mu\nu}$ definiert die Antwort der Materiewirkung auf Geometrieänderung. Der Proportionalitätsfaktor wird so gewählt, dass die Einstein-Gleichung die Standardform erhält.
\item Ein Erwartungswert ist eine Zahl pro Operator. Der Zustand enthält im Allgemeinen Erwartungswerte vieler Operatoren sowie Korrelationen. Verschiedene Zustände können denselben Erwartungswert für einen Operator besitzen.
\item Ein symmetrischer Tensor hat zehn Komponenten. Eichfreiheit und Nebenbedingungen entfernen unphysikalische Komponenten. Für ein masseloses Feld bleiben zwei Helizitäten.
\item Eichinvarianz verlangt erhaltene Quellen. Der universelle erhaltene Tensor der Materie ist der Energie-Impuls-Tensor.
\item Eine semiklassische Näherung nutzt Erwartungswerte auf klassischer Geometrie. Ein fundamentaler Abschluss muss die vollständige Quantenstruktur der gekoppelten Freiheitsgrade erfassen.
\item Setze $\ell=G^a\hbar^b c^c$ und vergleiche Dimensionen. Die Lösung ist $a=b=1/2$, $c=-3/2$, also $\ell_P=\sqrt{\hbar G/c^3}$.
\item Die Entropie ist proportional zur Fläche und enthält $\hbar$. Das deutet darauf, dass geometrische Horizonte mikroskopische quantenbezogene Zustände zählen.
\item Der Hauptsatz besagt: Der gemeinsame konsistente Abschluss von Quantenmaterie und dynamischer Raumzeit verlangt quantisierbare Gravitation.
\item Beispiele: nichtperturbative Definition, Messnachweis, Informationsproblem, Kandidatenauswahl.
\end{enumerate}

\chapter*{Schlusswort}
\addcontentsline{toc}{chapter}{Schlusswort}
Das dritte Semester schließt die Beweiskette. Der Weg führte von der Geometrie der Raumzeit über die Quantenstruktur der Materie zur linearen Gravitation, zum masselosen Spin-2-Feld, zur Selbstkopplung, zur semiklassischen Grenze, zur Planck-Skala, zum Pfadintegral, zur kanonischen Struktur und zu schwarzen Löchern.

Der Schluss ist klar:

\begin{merkkasten}
Wenn Materie quantenmechanisch ist, wenn Energie und Impuls die Quelle der Gravitation sind, wenn Raumzeit dynamisch ist, wenn Gravitation propagierende Freiheitsgrade besitzt und wenn der Gesamtprozess konsistent und unitär beschrieben werden soll, dann müssen gravitative Freiheitsgrade quantisierbar sein.
\end{merkkasten}

Das ist der mathematisch-physikalische Abschluss dieses Skripts. Die konkrete endgültige Theorie bleibt Aufgabe weiterer Forschung. Die Notwendigkeit des quantengravitativen Abschlusses folgt jedoch bereits aus den bekannten Prinzipien.

q.e.d.\\
Ingolf Lohmann

\end{document}
